\section{Introduction}
\label{sec:intro}

The rapid advancement of deep foundation models has transformed the paradigm of modern machine learning \cite{radford2021learning, brown2020language, devlin2018bert}. Rather than training monolithic systems from scratch, researchers and practitioners commonly fine-tune a pre-trained base model on individual tasks to produce specialized experts. However, when deployed in multi-task scenarios, hosting multiple large-scale networks simultaneously introduces severe computational, memory, and storage bottlenecks. Consequently, {\it model merging} has emerged as an active and vital area of research, aiming to consolidate the capabilities of multiple task-specific experts into a single unified network without the need for expensive joint retraining or multi-task fine-tuning \cite{ilharco2022editing, yadav2023ties, wortsman2022model}.

Early approaches to model merging focused on static, linear weight interpolation or task vector arithmetic \cite{ilharco2022editing, wortsman2022model}. While elegant and parameter-free, static merging methods often suffer from severe inter-task parameter interference and weight conflicts, causing performance degradation on complex tasks. To address this, recent works have shifted toward {\it test-time adaptive model merging}, where merging coefficients are optimized on unlabeled target test streams \cite{yang2024adamerging, jung2025symerge}. For example, AdaMerging \cite{yang2024adamerging} optimizes layer-wise coefficients, and SyMerge \cite{jung2025symerge} optimizes low-rank coordinate-scaling adapters. Pushing this trajectory to an extreme, FoldMerge \cite{foldmerge} constructs a non-linear weight-space diffeomorphism using a 4-layer RealNVP normalizing flow consisting of approximately 2.6 million parameters to warp and align expert parameters.

In this work, we take a step back and critically analyze this trajectory through the lens of Occam's razor. We argue that while modern test-time adaptive model merging achieves strong results, it is vital to map the boundary between parameter frugality and performance. We present a {\it critical, deconstructive audit} of test-time adaptive model merging, focusing on the role of parameter constraints under two under-discussed phenomena:
\begin{enumerate}
    \item {\bf Whole-Model Parameter Blending:} Restricting parameter-space blending to localized bottlenecks (such as a single layer) destroys fine-tuned task-vector capabilities, showing that multi-task adaptation requires whole-model coordinate interpolation.
    \item {\bf Parameter Constraint Boundaries:} Under extremely low-parameter regimes (e.g., our 8 task-wise scalars), weight-space optimization alone is too mathematically constrained to resolve parameter conflicts, making auxiliary classification head adaptation the primary driver of performance. In contrast, high-capacity layer-wise methods (like AdaMerging) can achieve substantial, genuine weight-space alignment gains due to their layer-wise degrees of freedom.
\end{enumerate}

To isolate these factors and evaluate adaptive merging under strictly controlled conditions, we introduce {\bf Barycentric Proximity-Anchored Merging (BPAM)}, a simple, low-parameter framework. BPAM reduces the weight-space blending parameters to the absolute minimum: exactly $K$ global task-wise scalars (where $K=8$ is the number of expert tasks). This represents an outstanding {\bf 99.99\% parameter footprint reduction} compared to FoldMerge's 2.6M parameters, and a {\bf 99.3\% parameter reduction} compared to AdaMerging's layer-wise parameters.

BPAM operates on three core minimalist principles:
\begin{itemize}
    \item {\bf Convex Barycentric Simplex Projection:} Rather than unconstrained parameter scaling (which scales up weight norms and distorts activations), we restrict our task merging coefficients to a convex barycentric simplex. This mathematically guarantees that the energy scale of the merged weights remains bounded and aligned with the original pre-trained representation space, preserving activation scales without complex normalizing flows.
    \item {\bf Mean-Field Proximity Regularization:} To stabilize the optimization landscape on small calibration streams, we introduce a soft closed-form $\ell_2$ penalty that pulls the individual task coefficients toward a uniform barycentric centroid. This simple geometric anchor acts as a powerful regularizer, preventing transductive overfitting without any extra computational or parameter cost.
    \item {\bf Teacher-Guided Adaptation:} We optimize our $K$ scalars using joint KL-divergence between the merged model predictions and individual expert teacher predictions on unlabeled target streams.
\end{itemize}

We evaluate BPAM on the standard 8-task image classification benchmark using CLIP ViT-B/32 across three distinct configurations:
\begin{itemize}
    \item {\bf BPAM-Restricted} (strictly 8 parameters on `model.visual.proj`, frozen heads) achieves an average accuracy of {\bf 51.38\%}, proving that projection-layer bottlenecking alone is insufficient and whole-model blending is necessary.
    \item {\bf BPAM-Static} (strictly 8 parameters on the whole image encoder, frozen heads) achieves {\bf 69.21\%} average accuracy, matching linear Task Arithmetic with only 8 parameters total and strictly frozen classifier probes.
    \item {\bf BPAM-Full} (8 task scalars + 388K classifier parameters) achieves {\bf 75.22\%} average accuracy, illustrating that for extremely parameter-pruned regimes, auxiliary classification head adaptation becomes the primary driver of performance.
\end{itemize}

By presenting this honest deconstruction, we establish a minimalist lower bound for adaptive merging, mapping the exact performance threshold where layer-wise degrees of freedom (as in AdaMerging) or high-capacity mapping models (as in FoldMerge and SyMerge) become indispensable. Our work provides a compelling boundary baseline that explores the trade-off between parameter frugality and multi-task weight blending, advancing the field towards mathematically-principled and well-understood optimization boundaries.

Our core contributions are summarized as follows:
\begin{itemize}
    \item We present a transparent, critical deconstructive audit exploring the physical boundaries of parameter-frugal test-time adaptive model merging.
    \item We introduce Barycentric Proximity-Anchored Merging (BPAM), a simple, low-parameter framework restricted to exactly $K$ scalars that serves as a conceptual baseline.
    \item We formulate a convex simplex constraint that preserves activation scales in a fully closed-form manner, and a Mean-Field Proximity Penalty designed to prevent transductive overfitting on test streams.
    \item We evaluate BPAM across three distinct parameter/head regimes, mapping the exact threshold where layer-wise and architectural degrees of freedom are needed for practical multi-task alignment.
\end{itemize}
